\chapter{結論與未來研究}\label{chap:conclusion}
\begingroup
\setlength{\emergencystretch}{2em}

\section{研究結論}

本研究以具有時間資訊的不平衡分類為對象，在美國企業年度財務資料上比較歷史與近期模型的運用方式。主要分析採用重建的下一年度破產事件標籤，並以公司最終狀態標籤作為對照。依三項研究問題，結論整理如下。

第一，近期資料的優勢依分類器而異。在相同年份配置與驗證規則下，XGBoost 與 LightGBM 以近期資料模型的 AUC 較高，Random Forest 以合併重訓最高，Logistic Regression 三種策略相近。因此，訓練資料期間應與分類器一起決定，不宜預設捨棄歷史資料或只用近期資料。

第二，在固定的 Old／New 模型池中，成員選擇比整合方式更重要。只整合近期模型的 New3 AUC 最高，依驗證資料選出兩個成員的靜態集成 F1-score 最高；等權平均全部六模型反而因納入較弱的歷史模型而降低 AUC，四種動態選擇方法也未超過上述兩者。實務上，保留單一期間與等權集成作為控制組，再以驗證資料決定是否加入歷史模型，比直接採用較複雜的動態方法更穩妥。

第三，特徵選取沒有帶來一致改善。保留 14 個特徵時大多接近完整特徵，保留 9 個特徵時 AUC 多半下降；Recall 與 F1-score 的方向則依集成方法而變。特徵選取的價值須連同精簡程度與漏判、誤報的權衡一起判斷。

此外，調參後的單一分類器與集成表現相當，而集成池成員使用預設參數。因此，本研究的證據支持「在固定模型池中審慎選擇成員」，尚不足以主張集成優於調參後的單一分類器。標籤定義也會改變結論：改用公司最終狀態標籤時，驗證選定的靜態集成 AUC 高於 New3，與事件標籤的結果相反。研究使用具時間資訊的資料時，應先確認標籤所代表的時間意義。

\section{研究限制}

本研究的限制如下。第一，事件標籤以破產公司的最後觀測年度推定，並非逐筆核對的破產申請日；驗證資料中的事件也只有 33 筆，選模與門檻決策的不確定性較大。第二，事件標籤結果尚未完成公司群聚配對檢定，主要差異屬描述性觀察。第三，集成池成員未經調參，Baseline 與集成之間的分數差同時受參數與驗證資料影響。第四，所有實驗採固定年份切分，測試期間不更新模型，因此不驗證持續學習或線上適應的能力。第五，實證只使用一份公開資料，同一公司也可能同時出現在訓練與測試期間，結論不延伸至未見公司或其他市場。

\section{未來研究方向}

後續研究可依上述限制延伸。首先，對事件標籤結果補作以公司為單位的配對檢定與多重比較校正，確認主要差異是否穩定。其次，以 Baseline 調參後的成員重建模型池，在相同驗證資料下比較集成與單一分類器，回答集成本身是否帶來增益。第三，採用滑動視窗或逐期更新的設計，檢驗歷史與近期模型在持續更新情境下的運用方式。第四，取得可核對的破產日期，或使用其他具有明確事件時間的資料，檢查結論能否推廣；並區分訓練期間出現過與未出現過的公司，評估對新個體的預測能力。最後，若要主張特徵選取可降低成本，應以相同程序量測訓練與推論時間。
\endgroup
